\section{Motivation}

Structured prediction tackles tasks whose output is a composite
object with interdependent components. Unlike standard
classification, where each input maps to a single label,
structured prediction must account for relationships between
output components. Typical examples include sequence labeling in
natural language processing, pixel-wise segmentation in computer
vision, and reasoning over combinatorial domains such as puzzles
and games.

Two aspects make structured prediction challenging. First, the
output space grows exponentially with the number of variables,
which renders exhaustive enumeration infeasible and calls for
specialized inference algorithms. Second, the training signal is
structured: a model must learn not only which local labels are
plausible, but also which joint configurations are globally
consistent. Both challenges are addressed by probabilistic
graphical models such as Markov Networks, which factorize the
score over local terms tied to the structure of the output.

In recent years, deep neural networks have become the dominant
approach for learning rich representations from raw inputs such
as images. Their predictions, however, are local and do not by
default enforce global constraints. A natural direction is
therefore to combine neural feature extractors with structured
output layers --- a neuro-symbolic architecture in which a neural
network proposes local evidence and a Markov Network aggregates
that evidence under combinatorial constraints.
SATNet~\cite{wang2019satnet} is one well-known instance of this
idea, replacing the Markov Network with a differentiable MAXSAT
solver.

This thesis investigates such a neuro-symbolic combination and
applies it to the task of solving Sudoku puzzles from both
symbolic and visual inputs. Sudoku is an appealing benchmark for
three reasons. Its constraint structure is fixed and
well-understood, so the contribution of the structured decision
layer can be isolated cleanly. Its solution is unique, so the
evaluation signal is unambiguous. And its visual variant
introduces a perceptual component (handwritten-digit recognition)
that must be learned jointly with the reasoning component,
making it a natural testbed for joint training of neural features
and structured output layers.

\section{Problem Statement}

The goal of this thesis is to design, implement, and evaluate a method
for training a Sudoku puzzle solver from examples of puzzle assignments
and their corresponding solutions. The solver must support two input
modalities: \emph{symbolic} grids, in which cells contain numerical
digits, and \emph{visual} grids, in which cells contain images of
handwritten digits, provided as 81 per-cell images.

Formally, the predictor maps an input $x$ to a configuration
$y = (y_1, \ldots, y_N)$ of cell values that is consistent with the
puzzle's row, column, and box constraints. The core of the predictor
is a neural network whose final decision layer is a Markov Network.
The method is exercised on two graph families:
\begin{itemize}
    \item \textbf{Chain structure.} The output is modeled as a linear
    chain. This variant is used for a Hidden-Markov-Chain (HMC)
    benchmark on which MAP inference is exact (Viterbi,
    $O(T K^{2})$), so the LP-relaxed learning loss (LP-M3N) can be
    compared head-to-head against exact max-margin training and the
    relaxation tightness can be measured empirically.
    \item \textbf{General-graph structure.} The output is modeled as a
    factor graph encoding Sudoku row, column, and box constraints. MAP
    inference is approximated by Belief Propagation.
\end{itemize}

Learning is cast in the max-margin structured-prediction framework. To
make training tractable on the cyclic Sudoku graph, the
structured-hinge loss is replaced by a linear-programming relaxation
following Franc \& Yermakov~\cite{franc2021} and
Franc, Pr\r{u}\v{s}a \& Yermakov~\cite{franc2022}. The relaxation is
derived in detail in Section~\ref{subsec:lp-m3n}.

The proposed system is compared against a two-stage baseline --- a
single-digit OCR classifier followed by a hard-coded Sudoku solver ---
using per-cell Hamming and per-puzzle $0/1$ losses
(Section~\ref{subsec:loss-functions}). This baseline separates
perception from reasoning and provides a reference point for measuring
the benefit of joint learning.

\section{Contributions of the Thesis}

The contributions of this thesis are:
\begin{enumerate}
    \item \textbf{Predictor design.} A neuro-symbolic predictor
    coupling a neural backbone with a Markov Network decision
    layer. The architecture handles both symbolic (per-cell
    one-hot) and visual (per-cell MNIST image) inputs, and
    supports both chain and general-graph output structures with
    no method-level changes between them.
    \item \textbf{Learning and inference.} An implementation of
    max-margin learning for the proposed predictor. Chain
    structures use exact dynamic programming (Viterbi) for
    loss-augmented inference; general graphs use the LP-relaxed
    M3N loss (LP-M3N), which avoids any inference oracle during
    training. At test time, decoding uses Viterbi on chains and
    Belief Propagation on general graphs.
    \item \textbf{Software library.} The Python library
    \texttt{mnlearn}, in which architecture, data paths,
    training hyperparameters, and the choice between
    structured-hinge and LP-M3N training are all specified via
    YAML. The library reuses standard PyTorch backbones and is
    organized so that every reported figure can be reproduced
    from a single \texttt{config.yaml} artifact written next to
    its run.
    \item \textbf{Benchmark.} A benchmark dataset covering two
    tasks (Sudoku and HMC) and two input modalities (symbolic
    and visual). Each task has $30{,}000\,/\,10{,}000\,/\,10{,}000$
    examples (train / val / test); visual variants are rendered
    from MNIST digits.
    \item \textbf{Empirical evaluation.} A statistical evaluation
    comparing the proposed predictor against the two-stage
    OCR + solver baseline using per-cell Hamming and per-puzzle
    $0/1$ losses. Learning curves are reported as a function of
    training-set size $N$, and a tightness check on the chain
    benchmark validates the LP-relaxed learning loss against
    exact structured-hinge training.
\end{enumerate}

\section{Structure of the Thesis}

The rest of the thesis is organized as follows:
\begin{itemize}
    \item \textbf{Chapter~\ref{chap:background} -- Theoretical
    Background.} Reviews structured prediction and the evaluation
    losses; Markov Networks and their factorization; MAP inference
    via Viterbi on chains and via LP relaxation on cyclic graphs;
    max-margin learning with the LP-relaxed surrogate; and neural
    feature extractors as unary providers.
    \item \textbf{Chapter~\ref{chap:datasets} -- Datasets.}
    Introduces the Sudoku and HMC datasets, detailing data
    generation, train/val/test splits, and dataset statistics for
    both symbolic and visual variants. The visual variants render
    MNIST digits from disjoint per-split pools to prevent image
    leakage.
    \item \textbf{Chapter~\ref{chap:method} -- Method.} Presents
    the proposed architecture (a neural backbone with a Markov
    Network decision layer), the two training paths
    (structured-hinge for chains, LP-relaxed M3N for general
    graphs), the evaluation-time inference decoders, the two-stage
    OCR baseline, and the configuration-driven Python library that
    implements them.
    \item \textbf{Chapter~\ref{chap:experiments} -- Experiments
    and Results.} Reports two experimental tracks: a tightness
    validation on the HMC chain benchmark (where LP-M3N is compared
    against exact structured-hinge training) and the main visual
    Sudoku results (compared against the two-stage OCR + solver
    baseline). Discusses findings, limitations, and failure modes.
    \item \textbf{Chapter~\ref{chap:conclusion} -- Conclusion.}
    Summarizes the contributions and outlines directions for future
    work.
\end{itemize}